\section{总结与延伸}

本讲建立了深度强化学习的整体框架。核心要点包括：

\begin{enumerate}
    \item 深度 RL 研究序贯决策问题及其解法，与监督学习有本质区别。
    \item RL 的应用已经非常广泛：从机器人到游戏到 LLM 的 post-training。
    \item 基本建模框架：状态、动作、策略、奖励、轨迹。
    \item 目标是最大化期望累计奖励 $\mathbb{E}_{\tau \sim p_\theta(\tau)} \left[\sum_t r(s_t, a_t)\right]$。
    \item 后续课程将依次介绍 imitation learning、policy gradients、actor-critic、Q-learning、offline RL、reward learning 和 RLHF。
\end{enumerate}

\subsection{拓展阅读}
\begin{itemize}
    \item Sutton \& Barto, \textit{Reinforcement Learning: An Introduction} (2nd edition)：经典 RL 教材。
    \item CS234 (Stanford)：更偏理论的 RL 课程，与 CS224R 互补。
    \item Sergey Levine, \textit{CS285 (Berkeley)}：另一门优秀的深度 RL 课程。
\end{itemize}

%% --- Fin del contenido --- %%

\end{document}
